% GENERATED FILE. Edit canonical lessons, then run npm run generate:books.
% canonical-source-hash: fnv1a64:52650afce463a47d
% canonical-lessons: LA-C361-argumentum, LA-C361-exemplum, LA-C361-documentum, LA-C361-demonstro, LA-C361-confirmo

\chapter{Proof and Example}
\label{ch:la-proof-and-example}
\hlchaptermodality{361}

\begin{chapteropening}
\textbf{By the end of this chapter:} \emph{I can bring an argument, an example and a proof, and say that something shows or confirms a point.}

Complete the last lesson of chapter 361: I can bring an argument, an example and a proof, and say that something shows or confirms a point.
\end{chapteropening}

\section[argūmentum, argūmentī]{argūmentum, argūmentī — an argument, a proof; a plot}

\label{lesson:LA-C361-argumentum}



\begin{quote}
Before the new one: say the Latin for foolishness, then the Latin for steadiness.
\end{quote}

\subsection*{You'll want to know: argūmentum, argūmentī}
\textbf{argūmentum, argūmentī} — \textquotedblleft{}an argument, a proof; a plot\textquotedblright{}.

\begin{quote}
Hoc argūmentum mihi nōn persuādet.
\end{quote}

\begin{quote}
\emph{This argument does not persuade me.}
\end{quote}

\emph{argūmentum, argūmentī} n. It is also the plot of a play. Built on \emph{arguō}, to make clear. English \emph{argument}.

\begin{cousinweb}
\emph{argūmentum} is built on \emph{arguō} \textquotedblleft{}to make clear, to prove\textquotedblright{} (EtymDB; Lewis \& Short), with the ending \emph{-mentum} that names a means. English \emph{argument} comes from it.
\end{cousinweb}

\begin{grammarlens}[title={What you've built}]
One more word for telling, explaining and describing at B1.
\end{grammarlens}

\subsection*{Your turn}
\begin{itemize}
\raggedright
  \item \emph{Say it:} \emph{argūmentum, argūmentī}
  \item \emph{Say it:} \emph{argūmentum, argūmentī}, once more
  \item \emph{Say it:} a sentence of your own with \emph{argūmentum, argūmentī}
  \item \emph{Recall:} say the Latin for foolishness, then the Latin for steadiness, then say \emph{argūmentum, argūmentī} again
\end{itemize}

\begin{tcolorbox}[breakable,colback=teal!4,colframe=teal!35!black,title={Before you move on}]
What does argūmentum, argūmentī mean? (\textquotedblleft{}An argument, a proof; a plot\textquotedblright{}.) Say it once more.
\end{tcolorbox}
\section[exemplum, exemplī]{exemplum, exemplī — an example; a warning}

\label{lesson:LA-C361-exemplum}



\begin{quote}
Before the new one: say the Latin for steadiness, then the Latin for an argument.
\end{quote}

\subsection*{You'll want to know: exemplum, exemplī}
\textbf{exemplum, exemplī} — \textquotedblleft{}an example; a warning\textquotedblright{}.

\begin{quote}
Pater exemplum bonum fīliō dedit.
\end{quote}

\begin{quote}
\emph{The father set his son a good example.}
\end{quote}

\emph{exemplum, exemplī} n. \emph{Exemplī grātiā}, for the sake of an example, gives English \emph{e.g.}; \emph{exemplī causā} means the same. \emph{Exemplum facere} is to make an example of someone.

\begin{cousinweb}
\emph{exemplum} comes from \emph{eximō} \textquotedblleft{}to take out\textquotedblright{} (EtymDB; Lewis \& Short), built from \emph{ex} and \emph{emō}, whose old sense was \textquotedblleft{}to take\textquotedblright{}: a piece taken out to show the rest. English \emph{example} and \emph{sample} come from it through French.
\end{cousinweb}

\begin{grammarlens}[title={What you've built}]
One more word for telling, explaining and describing at B1.
\end{grammarlens}

\subsection*{Your turn}
\begin{itemize}
\raggedright
  \item \emph{Say it:} \emph{exemplum, exemplī}
  \item \emph{Say it:} \emph{exemplum, exemplī}, once more
  \item \emph{Say it:} a sentence of your own with \emph{exemplum, exemplī}
  \item \emph{Recall:} say the Latin for steadiness, then the Latin for an argument, then say \emph{exemplum, exemplī} again
\end{itemize}

\begin{tcolorbox}[breakable,colback=teal!4,colframe=teal!35!black,title={Before you move on}]
What does exemplum, exemplī mean? (\textquotedblleft{}An example; a warning\textquotedblright{}.) Say it once more.
\end{tcolorbox}
\section[documentum, documentī]{documentum, documentī — a lesson, a proof; a warning example}

\label{lesson:LA-C361-documentum}



\begin{quote}
Before the new one: say the Latin for an argument, then the Latin for an example.
\end{quote}

\subsection*{You'll want to know: documentum, documentī}
\textbf{documentum, documentī} — \textquotedblleft{}a lesson, a proof; a warning example\textquotedblright{}.

\begin{quote}
Eius poena documentum omnibus fuit.
\end{quote}

\begin{quote}
\emph{His punishment was a lesson to everyone.}
\end{quote}

\emph{documentum, documentī} n. Built on the verb for teaching: something that teaches. In Latin it is a lesson or a proof, not a paper. English \emph{document}.

\begin{cousinweb}
\emph{documentum} is built on \emph{doceō} \textquotedblleft{}to teach\textquotedblright{} (EtymDB; Lewis \& Short), with the ending \emph{-mentum}: a thing that teaches. English \emph{document} comes from it.
\end{cousinweb}

\begin{grammarlens}[title={What you've built}]
One more word for telling, explaining and describing at B1.
\end{grammarlens}

\subsection*{Your turn}
\begin{itemize}
\raggedright
  \item \emph{Say it:} \emph{documentum, documentī}
  \item \emph{Say it:} \emph{documentum, documentī}, once more
  \item \emph{Say it:} a sentence of your own with \emph{documentum, documentī}
  \item \emph{Recall:} say the Latin for an argument, then the Latin for an example, then say \emph{documentum, documentī} again
\end{itemize}

\begin{tcolorbox}[breakable,colback=teal!4,colframe=teal!35!black,title={Before you move on}]
What does documentum, documentī mean? (\textquotedblleft{}A lesson, a proof; a warning example\textquotedblright{}.) Say it once more.
\end{tcolorbox}
\section[dēmōnstrō, dēmōnstrāre]{dēmōnstrō, dēmōnstrāre — to show, to point out; to prove}

\label{lesson:LA-C361-demonstro}



\begin{quote}
Before the new one: say the Latin for an example, then the Latin for a lesson.
\end{quote}

\subsection*{You'll want to know: dēmōnstrō, dēmōnstrāre}
\textbf{dēmōnstrō, dēmōnstrāre} — \textquotedblleft{}to show, to point out; to prove\textquotedblright{}.

\begin{quote}
Testis dēmōnstrāvit reum tum domī fuisse.
\end{quote}

\begin{quote}
\emph{The witness showed that the accused was at home then.}
\end{quote}

\emph{dēmōnstrō, dēmōnstrāre, dēmōnstrāvī, dēmōnstrātum}. With the accusative and infinitive it proves a claim. English \emph{demonstrate}.

\begin{cousinweb}
\emph{dēmōnstrō} is built from \emph{dē} and \emph{mōnstrō} \textquotedblleft{}to show\textquotedblright{} (Lewis \& Short), which belongs with \emph{moneō} \textquotedblleft{}to warn, to remind\textquotedblright{}. English \emph{demonstrate} comes from it.
\end{cousinweb}

\begin{grammarlens}[title={What you've built}]
One more word for telling, explaining and describing at B1.
\end{grammarlens}

\subsection*{Your turn}
\begin{itemize}
\raggedright
  \item \emph{Say it:} \emph{dēmōnstrō, dēmōnstrāre}
  \item \emph{Say it:} \emph{dēmōnstrō, dēmōnstrāre}, once more
  \item \emph{Say it:} a sentence of your own with \emph{dēmōnstrō, dēmōnstrāre}
  \item \emph{Recall:} say the Latin for an example, then the Latin for a lesson, then say \emph{dēmōnstrō, dēmōnstrāre} again
\end{itemize}

\begin{tcolorbox}[breakable,colback=teal!4,colframe=teal!35!black,title={Before you move on}]
What does dēmōnstrō, dēmōnstrāre mean? (\textquotedblleft{}To show, to point out; to prove\textquotedblright{}.) Say it once more.
\end{tcolorbox}
\section[cōnfirmō, cōnfirmāre]{cōnfirmō, cōnfirmāre — to strengthen; to confirm, to assert}

\label{lesson:LA-C361-confirmo}



\begin{quote}
Before the new one: say the Latin for a lesson, then the Latin for to show.
\end{quote}

\subsection*{You'll want to know: cōnfirmō, cōnfirmāre}
\textbf{cōnfirmō, cōnfirmāre} — \textquotedblleft{}to strengthen; to confirm, to assert\textquotedblright{}.

\begin{quote}
Litterae tuae opīniōnem meam cōnfirmant.
\end{quote}

\begin{quote}
\emph{Your letter confirms my opinion.}
\end{quote}

\emph{cōnfirmō, cōnfirmāre, cōnfirmāvī, cōnfirmātum}. Built on \emph{firmus}, firm: to make firm, then to state firmly. \emph{Animum cōnfirmāre}, to give someone courage. English \emph{confirm}.

\begin{cousinweb}
\emph{cōnfirmō} is built from \emph{con-} and a verb made on \emph{firmus} \textquotedblleft{}firm, strong\textquotedblright{} (EtymDB; Lewis \& Short): to make thoroughly firm. English \emph{confirm} comes from it.
\end{cousinweb}

\begin{grammarlens}[title={What you've built}]
One more word for telling, explaining and describing at B1.
\end{grammarlens}

\subsection*{Your turn}
\begin{itemize}
\raggedright
  \item \emph{Say it:} \emph{cōnfirmō, cōnfirmāre}
  \item \emph{Say it:} \emph{cōnfirmō, cōnfirmāre}, once more
  \item \emph{Say it:} a sentence of your own with \emph{cōnfirmō, cōnfirmāre}
  \item \emph{Recall:} say the Latin for a lesson, then the Latin for to show, then say \emph{cōnfirmō, cōnfirmāre} again
\end{itemize}

\begin{tcolorbox}[breakable,colback=teal!4,colframe=teal!35!black,title={Before you move on}]
What does cōnfirmō, cōnfirmāre mean? (\textquotedblleft{}To strengthen; to confirm, to assert\textquotedblright{}.) Say it once more.
\end{tcolorbox}
